\section{Preliminaries And The Relaxed Candidate Weight}

The sieve vocabulary used throughout:

\begin{itemize}
\item A \textbf{$P_2$ number} is a positive integer $n$ with
  $\Omega(n)\le2$, where $\Omega$ counts prime factors with multiplicity.
\item A \textbf{lower-bound sieve} is a sieve method that produces a lower bound for
  the count of elements of a sequence surviving a set of divisibility
  conditions; the quantity $\sum_{Q\le n\lt Q^2}a_Q(n)$ of this article is
  the sifted sum such a method would bound from below.
\item A \textbf{Type-I estimate} controls divisor sums: it bounds
  $\sum_d \tau(d)\,\bigl|A_d - A_1/\varphi(d)\bigr|$ for divisors $d$ up to
  some level $D$, i.e. it uses distribution of the sequence on arithmetic
  progressions. For this article the relevant remainders are the
  prime-progression discrepancies of Section~5.
\item A \textbf{Type-II estimate} controls bilinear sums
  $\sum_{m,n}\xi_m\kappa_n\,a_{mn}$ with arbitrary bounded coefficients,
  i.e. distribution of the sequence along products; Section~6 identifies the exact
  character modes such an estimate must control for this weight.
\end{itemize}

Choose a fixed exponent

\begin{equation*}
\frac13\lt\alpha\lt\frac12
\end{equation*}

and define

\begin{equation*}
\begin{aligned}
z&=X^\alpha=Q^{2\alpha},
\qquad
Z=P(z).
\end{aligned}
\end{equation*}

The upper restriction $\alpha\lt 1/2$ is useful because it gives $z\lt Q$, hence
$Z\mid W$. Define

\begin{equation*}
a_Q(n)
=
\mathbf1_{\gcd(n,W)=1}
\mathbf1_{\gcd(n+2,Z)=1}.
\end{equation*}

The candidate asks whether, for some fixed $\alpha>1/3$ and infinitely many
heads $Q$,

\begin{equation*}
\sum_{Q\le n\lt Q^2}a_Q(n)>0.
\end{equation*}

This is weaker than twin-prime positivity because $n+2$ is not required to
survive every prime below $Q$.
